\section{Framework and background.}\label{sec:framework}
\subsection{Notation}

We denote by $\sZ_- := \{\dots, -2, -1, 0\}$ the set of non-positive integers, by $\sN$ the non-negative integers, and for $m\in\sN^\ast$ we write $[m]:=\{1,\dots,m\}$; for any set $\gA$, $|\gA|$ denotes its cardinality. Random objects are written in uppercase (e.g.\ $X_t,Y_t$) and processes in bold roman (e.g.\ $\rmX=(X_t:t\in\sZ_-)$, $\rmIO=((X_t,Y_t):t\in\sZ_-)$), while a (finite) observed trajectory is written in bold lowercase, $\rvx=(x_t:t\in\sZ_-)\in\gIZ$ with $x_t\in\gI\subset\sR^d$ and $y_t\in\gY\subset\sR$; a length-$w$ window ending at time $\tau$ is $X_\tau^w:=(X_{\tau-w+1},\dots,X_\tau)\in(\gI)^w$ with realization $\rvx_\tau^w=(x_{\tau-w+1},\dots,x_\tau)$, and a windows dataset is $\gD=\{(\rvx_i,y_i)\}_{i=1}^N$ (we use $\E_P[\cdot]$ and $\Pr(\cdot)$ for expectation/probability under the law $P$). For an $n$-qubit register, $\gH:=\sC^{2^n}$, $\gSH$ is the space of density operators, and $\gBH$ is the space of bounded operators on $\gH$; we use Dirac notation, set $\ketbra{+}^{\otimes n} := \otimes_{i=1}^n \ketbra{+}$, and write $\braket{O}_{\rho}:=\Trc{O\rho}$ for the expectation of an observable $O$. We write $\normeuc{\cdot}$ for the Euclidean norm, and $\normHS{\cdot}$, $\normTr{\cdot}$, and $\normop{\cdot}$ for the Hilbert--Schmidt, trace, and operator norms, respectively. For spatial multiplexing with $R\ge1$ sub-reservoirs we work on $\gH^{\otimes R}:=\bigotimes_{r=1}^R\gH$ and $\gBH^R:=\gB(\gH^{\otimes R})$; with a mild abuse of notation, $\gSH^R$ denotes the corresponding state space and, in particular, the product states $\bigotimes_{r=1}^R\rho^{(r)}$ with each $\rho^{(r)}\in\gSH$. Finally, $\gO$ denotes the set of measured $k$-local observables, $m:\gSH^R\to\sR^{R|\gO|}$ the moment map producing features $\Phi=m\circ H^T$, and $\kappa$ is the Mat\'ern kernel with RKHS $\gHkp$ and norm $\normkp{\cdot}$; $K$ denotes the associated Gram matrix and $\lambda_{\rm reg}$ the Tikhonov regularization parameter.

%\TBDcomment{[State the difference between time series and processes]}

\subsection{Learning temporal data processes.}\label{subsec:learning-temporal-data}
Let 
$\rmIO = \pth{(X_t, Y_t): t \in \sZ_-}$, where $ X_t \in \gI$ and $Y_t \in \gY$, be a semi-infinite input/output stochastic process\footnote{The probability space $(\Omega, \gA, \sP)$ is fixed for all random variables and processes.}, with $\gI \subset \sR^{d}$ being some bounded measurable input set such that $\sup_{x \in \gI}\normeuc{x} \leq \Upsilon_{\scriptscriptstyle X}$, and $\gY \subset \sR$ some bounded measurable real-valued output set satisfying $\sup_{y \in \gY} |y| \leq \Upsilon_{\scriptscriptstyle Y}$.
We distinguish between the stochastic process and its realizations: the input/output process $\mathrm{IO}$ is a random object defined on a fixed probability space, while a (finite) time series corresponds to a single observed trajectory of this process, from which we construct windows for learning.
We consider the supervised learning problem where the task is to learn the data generating process (DGP) functional (map) $H^\star$ that assigns $Y_0$ to the process $\rmX = (X_t : t \in \sZ_-)$: $Y_0 = H^\star(\rmX)$, where the process $\rmIO$ is distributed according to some unknown distribution $P$~\cite{Gonon_Grigoryeva_Ortega_2020, qrc_risk_bounds_2025}. This problem is rooted in real-life processes (e.g., weather data) where we only have access to data up to date (with the most recent index being $t=0$) and we'd like to make a prediction $Y_0$ from such historical data (e.g., tomorrow's temperature, after a one-step shift of the time index).

As is central to machine learning, approximating the unknown DGP functional $H^\star$ follows from a risk minimization procedure. The (statistical) risk, or generalization error, associated with a functional $H$ is defined with respect to a fixed $L$-Lipschitz ($L > 0$) loss function $\ell: \sR \times \sR \longto \sR$ (e.g.\ squared error) as
\begin{equation}\label{eq:risk-def}
    R(H) := \E_P[\ell(H(\rmX), Y_0)].
\end{equation}
We are interested in learning $H^\star$ from the family $\gF$ of functionals satisfying the fading memory property (FMP)~\cite{Monzani_Prati_2025}. Intuitively, this property implies that if two time series $\rvx, \rvx' \in \gIZ$ are similar in their recent past (i.e., have similar values $x_t$ and $x'_t$ for times $t > t^\ast$ for some past time point $t^\ast \in \sZ_-$), then their outputs $H^\star(\rvx)$ and $H^\star(\rvx')$ will be close, even if $\rvx$ is very different from $\rvx'$ in the distant past ($t < t^\ast$).


Given a hypothesis class of functionals $\gC \subset \gF$, the ultimate goal of the learning procedure consists in determining the functional $H_\gC$ that exhibits minimal risk by solving
\begin{equation}\label{eq:in-class-optimization}
    H_\gC = \argmin_{H \in \gC} R(H).
\end{equation}
It is worth mentioning that the broader the class $\gC$ is, the more accurate the learning can be, hence the need for a rich (expressive) hypothesis class. Since the risk $R(H)$ depends on an unknown distribution $P$ as well as on a semi-finite time series process $\rmX$, it is generally infeasible to solve the optimization in~\cref{eq:in-class-optimization}. To this end, one must come up with an empirical risk $\hat{R}$ that does not deviate significantly from the true risk\footnote{Having a controllable $|R(H) - \hat{R}(H)|$ is what we mean by a good generalization.}, i.e., a quantity $|R(H) - \hat{R}(H)|$ that is controllable and can be computed efficiently. The learning procedure then becomes minimization of the empirical risk, i.e.\ $\argmin_{H \in \gC} \hat{R}(H)$.


\subsection{Quantum reservoir computers.}\label{subsec:qrc-background}

Crucial to reservoir computers (RCs)~\cite{Jaeger_2001_EchoState_GMD148} is the state-space evolution described as 
\begin{equation}\label{eq:RC-evolution}
    \rho_t = T(\rho_{t-1}, x_t), \quad \text{for all} \; t \in \sZ_-,
\end{equation}
where $T: \gS \times \gI \longto \gS$ is called the evolution map, and $\gS \subset \gB$ is the state space of some normed space $(\gB, \norm{\cdot})$. Quantum reservoir computers (QRCs) are an important class of RCs that leverage the state spaces of quantum systems~\cite{Chen_Nurdin_Yamamoto_2020}. Indeed, QRCs have as state space the set of quantum states described by the convex subset
\begin{equation}
    \gSH = \Set{\rho \in \gBH : \Trc{\rho} = 1, \rho \succeq 0},
\end{equation}
where $\gBH$ is the space of bounded linear operators on the Hilbert space $\gH = \sC^{2^n}$ of an $n$-qubit quantum system. Unless otherwise stated, we mainly consider the Hilbert--Schmidt (HS) norm (i.e., the Schatten-2 norm) $\normHS{\Psi} = \sqrt{\Trc{\Psi^\dag \Psi}}, \Psi \in \gBH$ as the norm on $\gBH$. Furthermore, at each time step, the state of a QRC is represented by a quantum state, while the quantum dynamics of the reservoir are described by a time-evolution map on the same quantum system (see \cite{Hayashi_Ishizaka_Kawachi_Kimura_Ogawa_2015} for details). Time-evolution maps on a quantum system are linear maps that are completely positive and trace-preserving (CPTP). For each fixed input $x_t \in \sR^{d}$, the induced map $T(\cdot, x_t): \gSH \longto \gSH$ is CPTP, and the state evolution is given by $\rho_t = T(\rho_{t-1}, x_t)$.


A QRC induces an inner filter, an inner-reservoir functional, and an outer-reservoir functional, defined respectively as
\begin{equation} \begin{aligned} 
    &U^T: \gIZ \longto {(\gSH)}^{Z_-}, \quad \;\; \rvx \longmapsto \big(T(s_{t-1}, x_{t}): t \in \sZ_- \big)\\ 
    &H^T: \gIZ \longto (\gSH), \qquad \;\; \rvx \longmapsto U^T(\rvx)_0\\ &H_h^T: \gIZ \longto \sR, \qquad\quad\quad\quad \; \rvx \longmapsto h \circ H^T(\rvx), 
\end{aligned} 
\end{equation}
where $h: \gSH \to \sR$ is some readout function that combines a measurement scheme conducted on the quantum system with some classical post-processing.

There are two properties that are crucial in reservoir computing: the echo state property (ESP) and the fading memory property (FMP). The first, ESP, is satisfied when for every sequence $\rvx \in \gIZ$ there exists a sequence of states $(\rho_t: t \in \sZ_-)$ that satisfies the relation in~\cref{eq:RC-evolution} for each $t \in \sZ_-$, and such a solution is unique. The ESP has attracted a lot of attention in the literature on reservoir computing~\cite{Jaeger_2001_EchoState_GMD148} because it gives the system a function-like definition where, for each input sequence, there is a single output sequence. The second property, on the other hand, states that a reservoir computer satisfies the FMP when its functional $H^T_h$ satisfies it in the sense described in the second paragraph of Section~\ref{subsec:learning-temporal-data}. Combined together, these properties imply that a state $\rho_t$ at time $t$ depends less and less on earlier states $\rho_{t-w}$ as we increase $w$. We will leverage this to show that the reservoir's initial state becomes increasingly irrelevant to the output as the reservoir processes more and more time points.


An interesting family of reservoir computers is that of evolution maps $T$ that are strictly contractive in their first argument:
\begin{equation}\label{eq:contraction}
%    \norm{T(\rho, x) - T(\rho', x)} \leq \lambda \norm{s - s'}, \; \text{for all} \; \rho, \rho' \in \gS, 
    \norm{T(\rho, x) - T(\rho', x)} \leq \lambda \norm{\rho - \rho'}, \; \text{for all} \; \rho, \rho' \in \gS,
\end{equation}
where $\lambda \in (0,1)$. This is of interest because any reservoir of this form automatically satisfies both the ESP and the FMP~\cite{Monzani_Prati_2025, Grigoryeva_Ortega_2018}. Subsequently, in Section~\ref{subsec:QR-embedding}, we propose a design of quantum reservoirs satisfying this condition, which leads to a principled hypothesis class $\gC_R$ of functionals that we leverage to approximate an unknown DGP functional $H^\star \in \gF$.